\cotefile{148}
\scanpage{44}
\mspage{1}

\uline{Groupoïdes de Teichmüller}.

\textcircled{1} Cas topologique.

La variante la plus simple est
\[
{\unsure{\mathscr T|\mathbb T}}_{\mathrm{top}}^{+}\quad\text{ou}\quad\mathscr T_{\mathrm{top}}
\]
(\uline{groupoïde de Teichmüller \uline{orienté} topologique})

% slanted
\margin{Finalement, tant qu'on ne travaille pas au point de vue arithmétique, il semble que c'est le gros $\mathscr T^{\pm}$ qui sera le plus utilisé, \deleted{\unsure{tout|tant}} c'est lui qui mérite donc ici la notation la plus simple $\mathscr T$, et le \unsure{gr. or|groupe|gr.} orienté de s'appeler $\mathscr T^{+}$.}

\uline{Objets}: surfaces \uline{à bord} compactes orientées

\uline{flèches}: classes d'isotopie d'Iso conservant l'orientation)

Variante immédiate
\[
\mathscr T_{\mathrm{top}}^{\pm}
\]
(\uline{groupoïde de Teichmüller} \uuline{biorienté} \uline{topologique}): mêmes objets, flèches $=$ \deleted{les} classes d'isotopie d'homéomorphismes (ne conservant pas néc. l'orientation).

Soit $[\pm1]$ le groupoïde ayant un seul objet $e$ (par exemple $e=\varnothing$) et comme groupe d'autom. $\pm1$. On a un morphisme de signature
\[
\mathscr T_{\mathrm{top}}^{\pm}\xrightarrow{\unsure{\chi}}\,[\pm1]
\]
dont la \uline{fibre} \deleted{est} (ou noyau) est $\mathscr T_{\mathrm{top}}^{+}=\mathscr T_{\mathrm{top}}$.

% slanted
\margin{\unsure{canulé|calculé}, si on \unsure{ne} se borne à des surfaces connexes !}

La variante la plus complexe est
\[
S_{*}\mathscr T{\unsure{\mathbb D|D}}_{\mathrm{top}}^{+}\quad\text{ou}\quad S_{*}\mathscr T\mathbb D_{\mathrm{top}},
\]
``\uline{groupoïde de Teichmüller} \inserted{\uline{orienté}} \uline{spécial dévissé topologique}''
% archive sheet number: 43

